% =====================================================================
%  Выпускная квалификационная работа
%  Прогнозирование направления цены биткоина: сравнение трёх постановок
%  3-классовой задачи классификации
%  Автор: Григорян Тигран Артурович
%  Научный руководитель: Слаболицкий Илья Сергеевич
%
%  Оформление воспроизводит стиль образцов (НИУ ВШЭ, ФКН):
%  титульный лист, двуязычная аннотация, нумерованные главы,
%  числовые ссылки [n], «Список литературы».
%  Компиляция: pdflatex (или xelatex) с поддержкой кириллицы.
% =====================================================================
\documentclass[12pt,a4paper]{report}

% Кириллица для обоих движков: pdflatex (T2A) и xelatex/lualatex (fontspec).
\usepackage{iftex}
\ifPDFTeX
  \usepackage[T2A]{fontenc}
  \usepackage[utf8]{inputenc}
\else
  \usepackage{fontspec}
  \defaultfontfeatures{Ligatures=TeX}
  \IfFontExistsTF{Latin Modern Roman}{\setmainfont{Latin Modern Roman}}{}
\fi
\usepackage[english,russian]{babel}
\usepackage[a4paper,left=30mm,right=15mm,top=20mm,bottom=20mm]{geometry}
\usepackage{amsmath,amssymb}
\usepackage{enumitem}
\usepackage{setspace}
\usepackage{indentfirst}
\usepackage[hidelinks]{hyperref}
\usepackage{titlesec}

\onehalfspacing
\setlength{\parindent}{1.25cm}

% Заголовки глав/разделов
\titleformat{\chapter}[block]{\normalfont\Large\bfseries\centering}{\thechapter}{1em}{}
\titlespacing*{\chapter}{0pt}{0pt}{20pt}

\begin{document}

% =====================================================================
%  ТИТУЛЬНЫЙ ЛИСТ  (формат — как у образцов; конкретные реквизиты — [заполнить])
% =====================================================================

\begin{titlepage}
    \centering
    \onehalfspacing

    {\small Федеральное государственное автономное образовательное учреждение высшего образования\par}
    \vspace{1em}
    {\bfseries\MakeUppercase{НАЦИОНАЛЬНЫЙ ИССЛЕДОВАТЕЛЬСКИЙ УНИВЕРСИТЕТ}\par}
    {\bfseries\MakeUppercase{«ВЫСШАЯ ШКОЛА ЭКОНОМИКИ»}\par}
    \vspace{1em}
    {\small Факультет компьютерных наук\par}
    {\small образовательная программа «Науки о данных»\par}

    \vfill

    {\Large\bfseries КУРСОВАЯ РАБОТА\par}
    \vspace{2em}

    %–– ЗАМЕНИТЕ текст внутри makebox на своё реальное название ––%
\begin{center}
 \bfseries            % жирный крупный шрифт
  % текст темы в "коробке", чтобы он переносился
  \parbox{0.9\textwidth}{\centering
    Предсказание высоковолатильных временных рядов с помощью моделей нейронных сетей}%   ← ваша тема
  \\[0.8ex]                  % чуть-чуть вертикального отступа
  \rule{0.9\textwidth}{0.4pt}% собственно линия (0.4 pt толщиной)
\end{center}
\vspace{1em}


    \vfill

    \begin{flushright}
    Выполнил:\\
    Студент группы \underline{\makebox[4cm][l]{группы МНКДН251}}\\[0.5em]
    \underline{\makebox[7cm][c]{Григорян Тигран Артурович}}\\[2em]

    Руководитель:\\[0.5em]
    \underline{\makebox[5cm][l]{преподаватель, магистр}}\\
    \underline{\makebox[7cm][c]{Слаболицкий Илья Сергеевич}}

    \end{flushright}

    \vfill

    {\centering Москва 2026\par}

\end{titlepage}
\tableofcontents


% =====================================================================
%  АННОТАЦИЯ / ANNOTATION
% =====================================================================
\chapter*{Аннотация}
\addcontentsline{toc}{chapter}{Аннотация}

Работа посвящена предсказанию высоковолатильных финансовых временных рядов с
помощью моделей нейронных сетей. В качестве канонического представителя такого
ряда выбран биткоин: задача формулируется как прогнозирование \emph{направления}
его цены на горизонте 24~часа~--- классификация на три класса $\{-1,0,+1\}$ на
часовых данных BTC/USDT с фиксированным порогом~$\pm 2\%$ (класс $+1$~--- значимый
рост, $-1$~--- значимое падение, $0$~--- отсутствие значимого движения, «флэт»).
Переход от прогноза уровня цены к прогнозу направления мотивирован тем, что
регрессия на уровни высоковолатильного ряда вырождается в ловушку случайного
блуждания и не несёт реальной предсказательной силы о движении.

Центральный объект изучения~--- \emph{нейросетевые модели} двух типов
представления данных: табличный многослойный перцептрон на инженерных признаках и
\emph{последовательностные} модели (GRU, LSTM, причинная свёрточная сеть~TCN) на
строго каузальных окнах исходных каналов. Исследуются три вопроса: (i)~способны ли
нейронные сети предсказывать направление высоковолатильного ряда и какие
архитектуры предпочтительны; (ii)~как способ постановки задачи влияет на
качество~--- сравниваются три декомпозиции одной и той же трёхклассовой задачи
(единая softmax-голова; две независимые бинарные модели; каскад
«гейт~$\to$~направление»); (iii)~даёт ли выразительная сила нейронных сетей
преимущество над \emph{сильными базовыми моделями} градиентного бустинга при
дефиците эффективно-независимых наблюдений, характерном для высоковолатильных
режимов.

Методология подчинена приоритету честности оценки: строго каузальные признаки и
каузальные окна, \mbox{purged}~$+$~embargoed walk-forward валидация как операционная
петля, отбор по коэффициенту Мэтьюса (MCC) вместо доли верных ответов,
экономическая оценка с издержками и статистические поправки на множественное
тестирование (Deflated Sharpe Ratio, Probability of Backtest Overfitting, тест
Дибольда--Мариано). Текст содержит постановку задачи, обзор литературы,
методологию, описание данных и признаков и дизайн экспериментов; раздел
эмпирических результатов оставлен заготовкой и будет заполнен после полного
вычислительного прогона.

\paragraph{Ключевые слова:} высоковолатильные временные ряды, нейронные сети,
глубокое обучение, рекуррентные сети, LSTM, GRU, TCN, прогнозирование направления,
классификация временных рядов, биткоин, purged walk-forward, нестационарность,
коэффициент Мэтьюса.

\chapter*{Annotation}
\addcontentsline{toc}{chapter}{Annotation}

This thesis addresses the prediction of highly volatile financial time series
using neural network models. Bitcoin is taken as the canonical instance of such a
series: the task is formulated as predicting the \emph{direction} of its price over
a 24-hour horizon~--- a three-class classification $\{-1,0,+1\}$ on hourly BTC/USDT
data with a fixed $\pm 2\%$ threshold ($+1$ a significant rise, $-1$ a significant
fall, $0$ the absence of significant movement, ``flat''). Moving from level to
direction forecasting is motivated by the fact that regression on the levels of a
highly volatile series collapses into the random-walk trap.

The central object of study is a family of \emph{neural network models} over two
data representations: a tabular multilayer perceptron on engineered features and
\emph{sequence} models (GRU, LSTM, and a causal temporal convolutional network,
TCN) on strictly causal windows of the raw channels. Three questions are
investigated: (i) whether neural networks can predict the direction of a highly
volatile series and which architectures are preferable; (ii) how the problem
formulation affects quality~--- three decompositions of the same three-class task
are compared (a single softmax head; two independent binary models; a
gate$\to$direction cascade); and (iii) whether the expressive power of neural
networks yields an advantage over \emph{strong gradient-boosting baselines} given
the scarcity of effectively independent observations characteristic of
high-volatility regimes.

Validation relies on strictly causal features and windows, purged and embargoed
walk-forward treated as an operational loop, selection by the Matthews correlation
coefficient, cost-aware economic evaluation, and corrections for multiple testing
(Deflated Sharpe Ratio, Probability of Backtest Overfitting, Diebold--Mariano).
The empirical results section is left as a placeholder to be completed after the
full computational run.

\paragraph{Keywords:} highly volatile time series, neural networks, deep learning,
recurrent networks, LSTM, GRU, TCN, direction forecasting, time-series
classification, Bitcoin, purged walk-forward, non-stationarity, Matthews
correlation coefficient.

% =====================================================================
%  ВВЕДЕНИЕ
% =====================================================================
\chapter*{Введение}
\addcontentsline{toc}{chapter}{Введение}

\section*{Актуальность темы исследования}
\addcontentsline{toc}{section}{Актуальность темы исследования}

Высоковолатильные финансовые временные ряды~--- одна из наиболее трудных для
прогнозирования категорий данных. Их характерные свойства~--- экстремальная
волатильность с кластеризацией, тяжёлые хвосты распределения, структурные сдвиги и
смена режимов~--- нарушают предположения классических статистических моделей.
Каноническим представителем такого ряда служит биткоин: за полтора десятилетия с
момента публикации его протокола~\cite{nakamoto2008} он прошёл путь от нишевого
криптографического эксперимента до класса активов с капитализацией, превышающей
триллион долларов США, торгуется круглосуточно и без выходных и демонстрирует
спекулятивные пузыри и быстрые обвалы. Рост институционального интереса и появление
производных инструментов усилили спрос на количественные методы прогнозирования и
управления риском на таких рынках.

В последнее десятилетие методы \emph{глубокого обучения} преобразили
прогнозирование временных рядов: рекуррентные сети (LSTM), свёрточные и
основанные на механизме внимания (трансформерные) архитектуры стали стандартом и
активно применяются к криптовалютам и другим высоковолатильным
рядам~\cite{ji2019, livieris2020, kim2022, dashtaki2025, wu2024, labbaf2024}.
Способность нейронных сетей улавливать нелинейные и временные зависимости делает их
естественным кандидатом для прогнозирования \emph{направления} движения
высоковолатильного ряда. Вместе с тем уместность нейросетевых моделей в этой задаче
далеко не очевидна: рынок близок к информационно эффективному, знак приращений цены
ведёт себя как близкий к белому шуму процесс, а прогноз \emph{уровня} упирается в
ловушку случайного блуждания~\cite{fama1970, lo1988}. Это смещает интерес к
прогнозированию направления и одновременно ставит вопрос: даёт ли выразительная
сила нейронных сетей реальное преимущество там, где сигнал слаб, а данных
эффективно мало?

Актуальность настоящей работы определяется двумя обстоятельствами. Во-первых,
несмотря на обилие нейросетевых моделей прогнозирования высоковолатильных рядов, в
литературе доминируют работы, предлагающие \emph{одну} архитектуру и демонстрирующие
её превосходство на единичном разбиении данных; систематического, честного по
утечкам сравнения нейросетевых архитектур между собой, с сильными базовыми моделями
и при разных способах \emph{постановки} задачи практически не проводится.
Во-вторых, значительная доля заявленных результатов получена без надлежащего
контроля утечек данных, нестационарности и множественного тестирования, что в
условиях большого числа параметров нейронных сетей особенно опасно и подрывает
воспроизводимость~\cite{lopez2018}. Настоящая работа отвечает на оба вызова,
проводя честную по методологии оценку нейросетевых моделей для предсказания
направления высоковолатильного ряда.

\section*{Цель и задачи исследования}
\addcontentsline{toc}{section}{Цель и задачи исследования}

\textbf{Цель исследования}~--- исследовать применимость моделей нейронных сетей к
предсказанию направления высоковолатильных финансовых временных рядов (на примере
биткоина), методологически строго и честно по утечкам сравнить нейросетевые
архитектуры между собой, с сильными базовыми моделями и при разных способах
постановки задачи, и установить границы предсказуемости направления на часовом
горизонте.

Для достижения цели поставлены следующие \textbf{задачи}:
\begin{enumerate}[leftmargin=*]
  \item провести обзор литературы по нейросетевым методам прогнозирования
        временных рядов и высоковолатильных финансовых рынков, методологии
        валидации финансовых моделей и метрикам качества;
  \item формализовать задачу прогнозирования направления как трёхклассовую
        классификацию и обосновать ключевые проектные решения (классификация
        вместо регрессии, третий «флэт»-класс, сравнение трёх постановок);
  \item собрать и подготовить данные: часовой ряд BTC/USDT и крипто-нативные
        экзогенные источники, спроектировать строго каузальный набор признаков и
        каузальное оконное представление для последовательностных нейросетей;
  \item реализовать нейросетевые модели (табличный многослойный перцептрон и
        последовательностные GRU/LSTM/TCN) в трёх постановках задачи, а также
        базовые модели градиентного бустинга для честного сравнения;
  \item построить честную процедуру валидации без утечек данных
        (purged~$+$~embargoed walk-forward), включая автоматические тесты утечки
        для нейронных сетей, и протокол отбора моделей;
  \item спроектировать сравнение нейросетевых архитектур, постановок задачи и
        базовых моделей по статистическим и экономическим критериям с поправками на
        множественное тестирование;
  \item сформулировать выводы о границах применимости нейронных сетей и направления
        дальнейшего улучшения.
\end{enumerate}

\section*{Объект и предмет исследования}
\addcontentsline{toc}{section}{Объект и предмет исследования}

\textbf{Объект исследования}~--- высоковолатильные финансовые временные ряды;
канонический представитель~--- биткоин (часовой ценовой ряд BTC/USDT и
сопутствующие крипто-нативные данные: деривативы, потоки сделок, on-chain
показатели, рыночный сентимент).

\textbf{Предмет исследования}~--- нейросетевые модели прогнозирования направления
движения высоковолатильного ряда, способы постановки (декомпозиции) трёхклассовой
задачи, а также методология честной валидации и сравнения нейронных сетей с
базовыми моделями.

\section*{Методы исследования}
\addcontentsline{toc}{section}{Методы исследования}

В работе используются методы \emph{глубокого обучения}~--- полносвязные
(многослойный перцептрон), рекуррентные (GRU, LSTM) и причинные свёрточные (TCN)
нейронные сети,~--- а в качестве сильных базовых моделей~--- градиентный бустинг
над решающими деревьями. Применяются методы математической статистики и анализа
временных рядов (проверки стационарности, автокорреляционный анализ, анализ режимов
волатильности), методология валидации финансового машинного обучения по Лопесу де
Прадо (purged $k$-fold, комбинаторная очищенная кросс-валидация,
walk-forward)~\cite{lopez2018}, аппарат статистической оценки значимости (Deflated
Sharpe Ratio, вероятность переобучения бэктеста, тест
Дибольда--Мариано)~\cite{bailey2014, bailey2017, diebold1995} и экономическое
бэктестирование торговой стратегии с явной моделью
издержек.

\section*{Научная новизна}
\addcontentsline{toc}{section}{Научная новизна}

Научная новизна работы состоит в следующем.
\begin{enumerate}[leftmargin=*]
  \item Впервые проводится \emph{like-for-like} сравнение трёх постановок
        трёхклассовой задачи направления (совместная softmax-голова против двух
        независимых бинарных моделей против каскада «гейт~$\to$~направление»)
        \emph{для нейросетевых моделей} высоковолатильного ряда при строго
        зафиксированных прочих условиях, что позволяет изолировать влияние самого
        способа декомпозиции на качество прогноза и экономику.
  \item Проведено честное (leakage-honest) сопоставление \emph{табличных и
        последовательностных} нейросетевых архитектур (многослойного перцептрона
        против GRU/LSTM/TCN) между собой и с сильными базовыми моделями
        градиентного бустинга при дефиците эффективно-независимых наблюдений,
        характерном для высоковолатильных режимов,~--- с каузальным оконным
        представлением и автоматическими тестами утечки, адаптированными к нейронным
        сетям; дополнительно предложен регим-нормализованный вариант целевой
        переменной (vol-scaled), корректирующий зависимость дисбаланса классов от
        рыночного режима.
  \item Реализована честная по утечкам и нестационарности методология оценки:
        purged~$+$~embargoed walk-forward как операционная петля, строгий контроль
        каузальности признаков и окон с автоматическими тестами утечки, а также
        сквозной учёт множественного тестирования через Deflated Sharpe Ratio и
        вероятность переобучения бэктеста по всему фактическому гриду конфигураций.
\end{enumerate}

\section*{Практическая значимость}
\addcontentsline{toc}{section}{Практическая значимость}

Практическая значимость работы определяется тем, что она даёт реалистичную,
воспроизводимую оценку границ применимости нейросетевых моделей к предсказанию
направления высоковолатильных финансовых рядов: какие архитектуры (табличные или
последовательностные) и при какой постановке задачи оправданы и в каких условиях
они уступают простым базовым моделям. Предложенная методология и программная
реализация образуют инструментарий для честной оценки предсказуемости, позволяют
формировать реалистичные ожидания относительно достижимого качества и служат
основой для построения и риск-контроля торговых стратегий, а также методическим
ориентиром для последующих исследований в области прогнозирования высоковолатильных
рядов методами глубокого обучения.

\section*{Структура работы}
\addcontentsline{toc}{section}{Структура работы}

Работа состоит из введения, четырёх глав, заключения и списка литературы. В
первой главе приводится обзор литературы и теоретические основы. Во второй главе
формализуется постановка задачи и излагается методология (три постановки, типы
целевой переменной, схема валидации, метрики и экономическая оценка). Третья глава
описывает данные и проектирование признаков. Четвёртая глава содержит дизайн
экспериментов; раздел эмпирических результатов оставлен заготовкой. В заключении
подводятся итоги и намечаются направления дальнейшей работы.

% =====================================================================
%  ГЛАВА 1
% =====================================================================
\chapter{Обзор литературы и теоретические основы}
\label{ch:review}

\section{Прогнозирование финансовых временных рядов: эффективность рынка и случайное блуждание}

Теоретической отправной точкой при моделировании цен финансовых активов служит
гипотеза эффективного рынка, согласно которой цены полностью и мгновенно отражают
доступную информацию~\cite{fama1970}. В слабой форме эффективности информационное
множество ограничено прошлыми ценами и объёмами; из неё следует, что приращения
цены не предсказуемы из прошлых цен сверх равновесной премии за риск, а ряд
лог-доходностей $r_t=\ln(S_t/S_{t-1})$ с точностью до первого приближения ведёт
себя как последовательность некоррелированных случайных величин. Эмпирически это
проявляется в практически плоской автокорреляционной функции доходностей при
наличии выраженной автокорреляции их модулей и квадратов~--- эффекте кластеризации
волатильности, впервые систематически описанном Мандельбротом~\cite{mandelbrot1963}.
Классические тесты случайного блуждания~--- расширенный тест Дики--Фуллера,
$Q$-статистика Льюнга--Бокса и тест отношения дисперсий~\cite{lo1988}~--- как
правило, не отвергают гипотезу единичного корня для уровня лог-цены и
подтверждают близость доходностей к белому шуму по первому моменту.

Для биткоина эти стилизованные факты~--- тяжёлые хвосты распределения доходностей,
кластеризация волатильности, асимметрия (эффект рычага) и нестационарность с
структурными сдвигами~--- многократно подтверждены эмпирически. Их совокупность,
с одной стороны, обосновывает трактовку цены как реализации случайного процесса, а
с другой~--- объясняет, почему прогнозирование \emph{уровня} цены обманчиво:
тривиальный прогноз «завтра как сегодня» обеспечивает высокий коэффициент
детерминации, не неся содержательной информации о направлении движения. Это
наблюдение является ключевым мотивом перехода к постановке задачи в терминах
направления, принятой в настоящей работе.

\section{Глубокое обучение для прогнозирования временных рядов}

Современная литература по прогнозированию финансовых и, в частности, криптовалютных
временных рядов методами глубокого обучения обширна; в ней доминируют рекуррентные
(LSTM, GRU), свёрточные и основанные на механизме внимания (трансформерные)
архитектуры. Цзи, Ким и Им~\cite{ji2019} сравнивают
семейства глубоких моделей (полносвязные сети, LSTM, свёрточные и остаточные сети)
для прогнозирования цены биткоина и приходят к практически важному выводу: при
переходе от регрессии уровня к \emph{классификации} направления (рост/падение)
классификационные модели оказываются полезнее для алгоритмической торговли, чем
регрессионные. Ливиерис и соавторы~\cite{livieris2020} исследуют ансамбли глубоких
моделей (LSTM, двунаправленные LSTM, свёрточные слои) на часовых данных
криптовалют, рассматривая одновременно регрессионную и классификационную
постановки и подчёркивая важность анализа автокорреляции ошибок прогноза.

Отдельное направление образуют работы, обогащающие ценовые признаки внешними
источниками. Ким и соавторы~\cite{kim2022} предлагают модель на основе
самовнимания над несколькими LSTM-модулями (SAM-LSTM), использующую \emph{on-chain}
данные~--- характеристики сетевой активности блокчейна~--- наряду с ценовыми, и
применяют детекцию точек смены режима для сегментации ряда. Орту и
соавторы~\cite{ortu2022} рассматривают именно \emph{классификацию направления}
дневной цены биткоина и эфира, сопоставляя технические и социально-сетевые
индикаторы; они сообщают рост точности с диапазона 51--55\% (только технические
признаки) до 67--84\% при добавлении социальных индикаторов, что прямо указывает на
ценность экзогенной информации и одновременно на скромный потолок предсказуемости
по чисто техническим признакам. Дашаки и соавторы~\cite{dashtaki2025} развивают
мультимодальный трансформерный подход с перекрёстным вниманием, объединяющий
рыночные данные и доменно-специфический сентимент (на основе FinBERT) для прогноза
тренда биткоина.

Обзорные и сравнительные исследования трансформерных и рекуррентных архитектур для
криптовалют~\cite{wu2024, labbaf2024} систематизируют это поле. Ву и
соавторы~\cite{wu2024} проводят обзор и эмпирическую оценку глубоких моделей для
прогнозирования цен криптовалют, отмечая преимущество многомерных
(многопризнаковых) подходов над одномерными и устойчивую конкурентоспособность
свёрточно-рекуррентных архитектур. Лаббаф Ханики и Мантури~\cite{labbaf2024}
комбинируют трансформер (с эффективным механизмом внимания) и двунаправленную
LSTM с техническими индикаторами на часовом и дневном горизонтах.

Из обзора видно, что: (i)~доминируют глубокие нейросетевые архитектуры
(рекуррентные, свёрточные, трансформерные), оцениваемые преимущественно на
единичных разбиениях; (ii)~классификация направления признаётся более релевантной
торговле, чем регрессия уровня; (iii)~экзогенные признаки (on-chain, сентимент,
поток сделок) систематически улучшают качество. При этом честного по утечкам
сравнения нейросетевых архитектур между собой и с сильными базовыми моделями, при
разных \emph{постановках} задачи и при последовательном контроле множественного
тестирования в этих работах, как правило, не проводится; вопрос о том, оправдывает
ли выразительная сила нейронных сетей их склонность к переобучению на
высоковолатильных рядах с малым числом независимых наблюдений, остаётся открытым.

\section{Классификация направления против регрессии уровней}

Выбор между прогнозированием уровня и направления является не техническим, а
концептуальным. Регрессия на уровень цены наследует свойства интегрированного
процесса первого порядка: тривиальный персистентный прогноз достигает высокого
$R^2$, который, однако, отражает лишь автокорреляцию уровней, а не способность
предвидеть движение. Постановка в терминах направления (рост / падение / отсутствие
значимого движения) не сводится к копированию последнего значения и непосредственно
соответствует торговому решению. Результаты~\cite{ji2019, ortu2022} эмпирически
поддерживают этот выбор: классификационные постановки оказываются практически
полезнее. Введение третьего, нейтрального класса дополнительно отделяет вопрос
«есть ли значимое движение» от вопроса «куда именно», что важно при доминировании
интервалов без выраженного тренда.

\section{Признаки: технические индикаторы, on-chain, деривативы и сентимент}

Множество предикторов в литературе можно сгруппировать в четыре класса.
\emph{Технические индикаторы} (скользящие средние, индикаторы импульса и
волатильности) описывают динамику самой цены и объёма; их применение
повсеместно~\cite{ji2019, livieris2020, labbaf2024}, однако предсказуемость по
ним ограничена. \emph{On-chain}-показатели~--- активность сети, число транзакций,
характеристики предложения~--- специфичны для блокчейн-активов и, как показано
в~\cite{kim2022}, несут дополнительный сигнал. \emph{Деривативные}
характеристики~--- ставка фондирования бессрочных фьючерсов и базис
«фьючерс--спот»~--- отражают плечо, премию за риск и позиционирование участников.
\emph{Сентимент}~--- социально-сетевые индикаторы и индексы настроений~--- улучшает
качество классификации направления~\cite{ortu2022, dashtaki2025}. Существенно, что
любой признак, чтобы быть допустимым предиктором, должен быть \emph{каузальным}
(использовать только информацию, доступную к моменту прогноза); нарушение этого
условия~--- частый источник утечек данных.

\section{Нестационарность, смена режимов и стилизованные факты}

Цены биткоина прошли несколько отчётливых режимов~--- фазы взрывного роста,
затяжные просадки, периоды штиля,~--- что проявляется в нестационарности уровня и
лишь слабой стационарности доходностей. Кластеризация волатильности и тяжёлые
хвосты~\cite{mandelbrot1963} означают, что мгновенная волатильность сама является
случайным, изменяющимся во времени процессом. Практическое следствие для
прогнозирования направления двояко. Во-первых, модель, обученная в одном режиме,
систематически теряет качество в другом, что требует адаптивных схем обучения и
валидации. Во-вторых, поскольку предсказуема в первую очередь именно
\emph{волатильность} (а не знак доходности), потенциальный источник прогностической
силы~--- обусловливание на состоянии волатильности, режиме и внутрисуточной
сезонности, а не линейный прогноз доходности.

\section{Методология валидации финансового машинного обучения и переобучение}

Корректная оценка предсказательной модели на финансовых данных требует учёта
зависимости и многобаровых, перекрывающихся меток. Лопес де Прадо~\cite{lopez2018}
показывает, что обычная $k$-fold кросс-валидация переоценивает качество, поскольку
перекрытие меток и rolling-признаков создаёт утечку между обучающей и тестовой
выборками; в качестве решений предложены \emph{purged} $k$-fold (с очисткой
перекрывающихся меток и буфером-эмбарго) и \emph{комбинаторная очищенная
кросс-валидация} (CPCV), порождающая множество путей бэктеста. Для прогнозирования
с разворачивающимся временем естественной является walk-forward (forward-chaining)
схема, в которой обучающая выборка всегда строго предшествует тестовой; такую схему
можно трактовать как симуляцию эксплуатации модели.

Множественное тестирование~--- перебор большого числа конфигураций~--- порождает
риск отбора спурьёзно «хорошей» стратегии. Бэйли и Лопес де Прадо~\cite{bailey2014}
вводят \emph{Deflated Sharpe Ratio}, корректирующий оценку коэффициента Шарпа на
число испытаний и ненормальность доходностей; Бэйли и соавторы~\cite{bailey2017}
предлагают оценку \emph{вероятности переобучения бэктеста} (PBO) через
комбинаторно-симметричную кросс-валидацию. Тест Дибольда--Мариано~\cite{diebold1995}
позволяет статистически сравнивать предсказательную точность двух моделей. Эти
инструменты образуют каркас честной оценки, принятый в настоящей работе.

\section{Метрики качества: коэффициент Мэтьюса и экономические критерии}

При сильном дисбалансе классов доля верных ответов (accuracy) вводит в заблуждение:
тривиальный прогноз доминирующего класса достигает высокой точности, не неся
информации. Коэффициент корреляции Мэтьюса (MCC)~\cite{matthews1975}~--- корреляция
Пирсона на матрице ошибок~--- учитывает все её ячейки и в среднем информативнее
точности и $F_1$ при дисбалансе, поэтому используется как первичная метрика отбора.
Наряду с классификационными метриками решающее значение имеет \emph{экономическая}
оценка: перевод прогнозов в торговый сигнал и расчёт прибыльности с учётом издержек
(коэффициенты Шарпа, Сортино, Кальмара, максимальная просадка), а также
precision-at-coverage как торговой метрики качества решения. Окончательным арбитром
значимости служит статистический контроль переобучения (DSR, PBO), описанный выше.

\section{Выводы по главе и обоснование настоящей работы}

Проведённый обзор позволяет выделить пробел в литературе. Несмотря на обилие
нейросетевых моделей прогнозирования высоковолатильных рядов, (i)~отсутствует
честное по утечкам \emph{like-for-like} сравнение нейросетевых архитектур (табличных
и последовательностных) между собой и с сильными базовыми моделями при разных
способах постановки трёхклассовой задачи направления; (ii)~значимая часть
результатов получена без строгого контроля утечек, нестационарности и
множественного тестирования, что особенно опасно для моделей с большим числом
параметров; (iii)~зависимость дисбаланса классов от рыночного режима, как правило,
не учитывается в определении целевой переменной. Настоящая работа закрывает этот
пробел, оценивая нейросетевые модели в трёх постановках задачи на фоне сильных
базлайнов, с регим-нормализованным вариантом метки и сквозным учётом переобучения.
Формальная постановка задачи и методология излагаются в Главе~\ref{ch:method}.

% =====================================================================
%  ГЛАВА 2
% =====================================================================
\chapter{Постановка задачи и методология}
\label{ch:method}

\section{Формальная постановка трёхклассовой задачи}

Пусть $\{S_t\}$~--- часовой ряд торгуемой цены закрытия BTC/USDT, $H=24$~---
горизонт прогноза в часах, $\theta$~--- порог значимого движения. Доходность за
горизонт в момент~$t$ определяется как $\mathrm{ret}_t = S_{t+H}/S_t - 1$, а метка
$y_t\in\{-1,0,+1\}$~--- как
\begin{equation}
y_t =
\begin{cases}
\,+1, & \mathrm{ret}_t \ge +\theta,\\[2pt]
\,-1, & \mathrm{ret}_t \le -\theta,\\[2pt]
\,\phantom{+}0, & \text{иначе},
\end{cases}
\qquad \theta = 2\%.
\end{equation}
Содержательно модель отвечает не на вопрос «какой будет цена», а на вопрос «уйдёт
ли цена за~$\pm 2\%$ вверх или вниз в течение суток или останется в коридоре».
Порог~$\theta$ трактуется как \emph{определение целевой переменной}, а не как
настраиваемый гиперпараметр: его недопустимо подбирать под классификационную
метрику, поскольку при разных порогах изменяются базовые частоты классов и
сравнение метрик становится некорректным (форма скрытого подглядывания в разметку).

\paragraph{Почему классификация, а не регрессия.} Как показано в
Главе~\ref{ch:review}, регрессия на уровень цены вырождается в ловушку случайного
блуждания. Постановка в терминах направления обходит эту ловушку и ближе к
реальному торговому решению.

\paragraph{Почему три класса.} Третий, нейтральный класс «флэт» вводится порогом и
принципиально важен: большинство 24-часовых окон не сдвигаются на~$2\%$, и
требование направленной ставки там, где значимого движения нет, означает навязывание
сигнала в шуме. Флэт~--- это класс, в котором \emph{не торгуют}; такая постановка
честнее бинарной «вверх/вниз» и отделяет наличие движения от его направления.

\section{Три постановки задачи}

Все три постановки реализуют единый интерфейс: на входе~--- признаки (или каузальное
окно для последовательностных моделей), на выходе~--- распределение по трём классам
$[\,P(-1),P(0),P(+1)\,]$, прогноз~--- класс с максимальной вероятностью, а
направленный балл уверенности определяется как $\mathrm{score}=P(+1)-P(-1)$.
Постановка~--- это \emph{способ декомпозиции} задачи, не зависящий от конкретного
обучателя; чтобы изолировать его влияние, обучатель внутри одного сравнения
фиксируется. В центре внимания настоящей работы~--- нейросетевые обучатели
(раздел~\ref{sec:nn}); градиентный бустинг~\cite{ke2017} используется как сильная
базовая модель.

\paragraph{Постановка A (мультиклассовый softmax).} Одна мультиклассовая модель
напрямую оценивает $[\,P(-1),P(0),P(+1)\,]$.

\paragraph{Постановка B (две независимые бинарные модели).} Обучаются две модели~---
$p_{\text{up}}=\Pr[\mathrm{ret}\ge+\theta]$ и $p_{\text{dn}}=\Pr[\mathrm{ret}\le-\theta]$,
обе на всех данных; затем вероятности согласуются в три класса по
зафиксированному до теста правилу
\[
P(+1)\propto p_{\text{up}}(1-p_{\text{dn}}),\quad
P(-1)\propto (1-p_{\text{up}})p_{\text{dn}},\quad
P(0)\propto (1-p_{\text{up}})(1-p_{\text{dn}})+p_{\text{up}}p_{\text{dn}},
\]
с последующей нормировкой. Для каждой стороны (long/short) допускаются отдельные
пороги уверенности.

\paragraph{Постановка C (каскад «гейт~$\to$~направление»).} Сначала «гейт»
оценивает вероятность значимого движения $p_{\text{big}}=\Pr[\,|\mathrm{ret}|\ge\theta\,]$,
затем модель направления различает рост и падение; итоговые вероятности~---
$P(+1)=p_{\text{big}}\,p_{\text{up}}$, $P(-1)=p_{\text{big}}(1-p_{\text{up}})$,
$P(0)=1-p_{\text{big}}$. Рассматриваются два варианта обучения модели направления:
\emph{c1}~--- на реальной популяции крупных движений ($|\mathrm{ret}|\ge\theta$); и
\emph{c2}~--- на популяции, отобранной/взвешенной по предсказанию гейта (с
демонстрацией смещения отбора, selection bias). Чтобы вариант c2 оставался
корректной демонстрацией, веса гейта оцениваются вне обучающей подвыборки
(out-of-fold).

\textbf{Исследовательский вопрос} формулируется так: влияет ли сам способ
декомпозиции (совместная softmax-голова против независимых бинарных моделей против
каскада) на качество прогноза и экономику, если все прочие элементы конвейера
зафиксированы?

\section{Типы целевой переменной}

Помимо способа моделирования, рассматриваются три \emph{определения} целевой
переменной, каждое из которых является отдельным сквозным экспериментом и
сравнивается по экономике, а не по классификационным метрикам (различные базовые
частоты делают их несопоставимыми).
\begin{itemize}[leftmargin=*]
  \item \textbf{Pointwise (point-to-point).} Метка определяется по конечной точке:
        $\mathrm{ret}_t=S_{t+H}/S_t-1$ с фиксированным порогом~$\pm\theta$; правило
        выхода в бэктесте~--- удержание горизонта~$H$.
  \item \textbf{Triple-barrier.} Путь-зависимая разметка: от уровня $S_t$ строятся
        верхний ($+\theta$) и нижний ($-\theta$) барьеры и вертикальный барьер на
        горизонте~$H$; метка определяется тем, какой барьер достигнут \emph{первым}.
        Правило выхода~--- по касанию барьера (стоп/тейк), что соответствует
        реальной сделке со стоп-лоссом.
  \item \textbf{Vol-scaled (регим-нормализованный).} Порог адаптируется к текущей
        каузально оценённой волатильности: $\theta_t=k\,\sigma_t$, где $\sigma_t$~---
        экспоненциально сглаженная волатильность доходностей со сдвигом на один бар
        (без заглядывания вперёд), масштабированная множителем~$\sqrt{H}$; параметр
        $k$ перебирается по небольшому пред-зарегистрированному набору. Такой таргет
        корректирует зависимость дисбаланса классов от режима волатильности.
\end{itemize}
Порог метки согласуется с правилом выхода: point-to-point~--- удержание~$H$;
triple-barrier~--- стоп/тейк по барьеру.

\section{Нейросетевые модели и базовые модели сравнения}
\label{sec:nn}

Центральный объект изучения~--- \emph{нейросетевые модели} в двух представлениях
данных.

\paragraph{Табличный многослойный перцептрон (MLP).} Полносвязная сеть на тех же
инженерных признаках, что и базовые модели (apples-to-apples сравнение):
1--2~скрытых слоя, нормировка батча, dropout; на выходе~--- softmax по трём классам.
Предобработка (импутация пропусков медианой, стандартизация) обучается строго на
train-фолде.

\paragraph{Последовательностные модели (GRU, LSTM, причинный TCN).} Работают на
строго каузальных окнах $[t-L+1,\dots,t]$ исходных каналов (лог-доходности,
лог-объём, несколько нормированных индикаторов), организованных в тензор формы
$(L,C)$. Для свёрточной сети используется \emph{причинная} (causal) свёртка,
исключающая заглядывание вперёд; рекуррентные сети берут последнее скрытое
состояние. Это позволяет нейросети извлекать временные паттерны сверх
hand-engineered лагов.

\paragraph{Обучение и регуляризация.} Кросс-энтропия (или focal loss) с весами
семплов (класс~$\times$~уникальность~$\times$~давность); оптимизатор AdamW с
весовым затуханием; ранняя остановка по MCC на отложенном purged-валидационном
срезе; агрессивная регуляризация (малые сети, dropout, weight decay)~--- ввиду
дефицита эффективно-независимых наблюдений. Нейросетевая часть исполняется в
отдельном процессе (во избежание конфликтов рантаймов с базовыми библиотеками) и при
наличии~--- на графическом ускорителе.

\paragraph{Базовые модели (baselines).} В качестве сильных базовых моделей
используется градиентный бустинг над решающими деревьями~--- LightGBM~\cite{ke2017}
(основная), а также XGBoost~\cite{chen2016} и CatBoost~\cite{prokhorenkova2018}: они
эффективны на табличных данных указанного масштаба, нативно работают с пропусками и
детерминированы. Честное сопоставление нейросетей с этими базлайнами на одних и тех
же purged-фолдах является ядром работы. При этом, как отмечается в
литературе~\cite{ji2019, wu2024} и согласуется с оценкой эффективного числа
независимых наблюдений (непересекающихся 24-часовых меток), на табличных данных
такого масштаба сильные базлайны часто оказываются конкурентоспособными, а
нейросетевые модели склонны к переобучению; последовательностные архитектуры~---
наиболее перспективная для нейросетей ниша, особенно на больших объёмах данных.

\section{Валидация: purged + embargoed walk-forward как операционная петля}

Основным, deploy-честным методом отбора и оценки является \emph{purged~$+$~embargoed
walk-forward} (forward-chaining): обучающая выборка всегда строго предшествует
тестовой; перед каждым тестовым блоком из обучения удаляются бары, чьи 24-часовые
метки перекрываются с тестовым блоком (purge), и добавляется буфер-эмбарго,
устраняющий утечку через автокорреляцию. Строгий контроль границы требует
выполнения условия $\max(\text{train}) + H < \min(\text{test})$, исключающего
касание меткой последнего обучающего бара тестового периода. Такая схема трактуется
как \emph{операционная петля}: каждый шаг~--- «переобучились на окне~$\to$
предсказали следующий период~$\to$ сдвинулись вперёд», то есть петля валидации
совпадает с петлёй эксплуатации. Для учёта нестационарности предусмотрены
адаптационные настройки (расширяющееся или скользящее окно, взвешивание по
давности, периодичность переобучения), рассматриваемые как абляции, а не как
настраиваемые в оптимизации параметры.

В качестве \emph{вторичной} оценки робастности используется комбинаторная очищенная
кросс-валидация (CPCV), порождающая множество путей бэктеста и служащая основой для
оценки вероятности переобучения; при этом учитывается, что CPCV обучает в том числе
на пост-тестовых данных и потому является менее deploy-честной. Предобработка
(шкалирование, кодирование, импутация) выполняется строго внутри обучающего фолда,
чтобы исключить утечку.

\section{Метрики, статистическая значимость и экономическая оценка}

Обучение моделей ведётся по логарифмической функции потерь; \emph{отбор и отчёт}~---
по коэффициенту Мэтьюса (как первичной метрике, устойчивой к дисбалансу), а также
по macro-$F_1$, полноте по классам, матрице ошибок и калибровочным метрикам;
точность (accuracy) как целевая не используется. Подбор гиперпараметров проводится
методом байесовской оптимизации на purged walk-forward (не случайным поиском и не
обычным $k$-fold), с отбором по среднему MCC по фолдам.

Прибыльность оценивается \emph{экономическим бэктестом} с явными допущениями: сигнал
формируется на закрытии бара~$t$, вход осуществляется по цене открытия следующего
бара~$t{+}1$ (а не по тому же закрытию), удержание~--- на горизонт (или по барьеру
для triple-barrier); допускается короткая позиция, размер позиции фиксирован;
издержки задаются в базисных пунктах. Чтобы оценка не зависела от произвольной фазы
входа, метрики усредняются по всем фазовым сдвигам входа, а коэффициент Шарпа
считается на де-перекрытых по-сделочных доходностях с аннуализацией по числу
\emph{сделок}, а не баров. Обязательным является сравнение с пассивной стратегией
«купи и держи». Статистическая значимость контролируется через Probabilistic и
Deflated Sharpe Ratio~\cite{bailey2014}, вероятность переобучения
бэктеста~\cite{bailey2017} и тест Дибольда--Мариано~\cite{diebold1995} между
постановками, причём в число испытаний для DSR включается \emph{весь} фактически
перебранный грид конфигураций.

% =====================================================================
%  ГЛАВА 3
% =====================================================================
\chapter{Данные и признаки}
\label{ch:data}

\section{Базовый ценовой ряд}

В качестве базового ряда используются часовые свечи OHLCV пары BTC/USDT, полученные
с биржи Binance, за весь доступный непрерывный период~--- ориентировочно с августа
2017~г. по июнь 2026~г. (около $7.7\cdot 10^{4}$ баров, $\approx 8.8$ года). Выбор
обоснован следующим. Во-первых, это наиболее длинная непрерывная история по
биткоину, необходимая для устойчивой валидации по множеству рыночных режимов.
Во-вторых, используется именно \emph{торгуемая} цена (а не взвешенная комбинация
OHLC), поскольку на ней же строится экономический бэктест,~--- чтобы сигнал и
исполнение находились в единой системе координат. В-третьих, переиспользуется уже
имеющийся загрузчик данных без введения новой зависимости, что обеспечивает
консистентность слоя данных. Дата отсечки фиксируется заранее (а не привязывается к
текущему моменту) ради воспроизводимости; сырой снимок данных сохраняется вместе с
контрольной суммой и числом строк.

Пропуски в истории немногочисленны (порядка~$10^{2}$ баров, ранние сбои API);
они обрабатываются реиндексацией на полную часовую сетку с явной пометкой гэпа,
чтобы доходность не вычислялась «через дыру» за интервал, больший одного бара.

\section{Экзогенные данные: крипто-нативные, а не макроэкономические}

Экзогенный набор по умолчанию~--- крипто-нативный, обладающий длинной историей и не
требующий платных ключей доступа:
\begin{itemize}[leftmargin=*]
  \item \textbf{Деривативы:} ставка фондирования (funding) бессрочных фьючерсов и
        базис «фьючерс--спот»~--- индикаторы плеча, премии за риск и позиционирования;
  \item \textbf{Поток сделок (order-flow):} дисбаланс агрессивных покупок и продаж
        и его производные, вычисляемые из объёмов тейкер-покупок и доступные на всю
        историю;
  \item \textbf{On-chain:} дневные показатели сетевой активности (активные адреса,
        число транзакций, хешрейт, характеристики предложения), мотивированные
        результатами~\cite{kim2022};
  \item \textbf{Сентимент:} индекс «страха и жадности» (Fear~\&~Greed).
\end{itemize}

Кросс-активы с традиционных рынков (фондовые индексы, валюты, металлы) сознательно
\emph{не} включены в базовую модель: их часовая история ограничена примерно~729
днями, и их добавление схлопнуло бы обучающее окно для всех признаков примерно до двух лет,
пожертвовав длинной историей; к тому же на часовом масштабе они «протухают» вне
биржевых сессий, тогда как криптовалюта торгуется непрерывно. Поэтому кросс-активы
оставлены опциональной \emph{абляцией} на коротком окне. Открытый интерес (open
interest) исключён из длинных признаков ввиду крайне ограниченной бесплатной истории.

\section{Проектирование признаков}

Все признаки строго трейлинговые (каузальные)~--- используют только информацию,
доступную к моменту прогноза. Набор включает:
\begin{itemize}[leftmargin=*]
  \item логарифмические доходности и их лаги на наборе горизонтов
        $\{1,2,3,6,12,24,48,72,168\}$;
  \item rolling-волатильность на окнах $\{24,72,168\}$ и её отношения (терм-структуру
        волатильности);
  \item \emph{нормированные} технические индикаторы~--- гистограмму MACD (а не сырую
        линию, дрейфующую со шкалой цены), отношение ATR к цене, z-оценку изменения
        OBV;
  \item логарифм объёма и его изменения/z-оценки; скользящие средние в форме
        отношений $\mathrm{close}/\mathrm{MA}-1$;
  \item календарные признаки в форме синус/косинус (час, день недели, месяц);
  \item \textbf{режим/сезонные} признаки: диурнальную климатологию волатильности
        (типичную абсолютную доходность для данного сочетания «час~$\times$~день
        недели», вычисляемую каузально), ранг и z-оценку текущей волатильности
        относительно недавней истории, просадку от трейлинг-максимума, длинный
        тренд и 30-дневную волатильность.
\end{itemize}
Экзогенные признаки лагируются у источника (значение в строке~$t$ отражает состояние,
известное строго до~$t$); дневные on-chain показатели сдвигаются не менее чем на
публикационную задержку до приведения к часовой сетке. Метка \emph{не} кэшируется, а
строится «на лету» из закэшированных признаков и OHLC, что делает смену типа целевой
переменной мгновенной.

Для последовательностных нейросетей те же каузальные каналы организуются в скользящие
окна $[t-L+1,\dots,t]$ длины~$L$ (тензор формы $(L,C)$): окно содержит только прошлое,
поэтому не вносит утечки метки, а первые $L-1$ баров каждого фолда, не имеющие полного
окна, отбрасываются. Табличные модели (бустинг, MLP) используют те же признаки в
плоском виде.

\section{Пропуски, структурное отсутствие и воспроизводимость}

Часть экзогенных признаков существует не на всей истории: ставка фондирования и
базис появляются лишь с запуска бессрочных фьючерсов (сентябрь~2019~г.), а спот~---
с августа~2017~г.; отсюда значительная доля пропусков одним непрерывным куском в
ранней эпохе (аналогично индексу настроений с начала~2018~г.). Это \emph{структурное}
отсутствие информативно и потому не устраняется удалением: деревья принимают
пропуски нативно и кодируют «до-перп»-эпоху отдельной ветвью. Для моделей, не
допускающих пропусков (нейросети, логистическая регрессия), вводятся бинарные флаги
наличия соответствующих рядов, а влияние деривативных признаков оценивается
\emph{абляцией} (с/без funding и базиса), а не слепым удалением. Для всех признаков,
содержащих rolling- и сезонные статистики, проводится отдельная проверка
каузальности.

% =====================================================================
%  ГЛАВА 4
% =====================================================================
\chapter{Дизайн экспериментов}
\label{ch:design}

\section{Принципиальный эксперимент: нейросети против базлайнов}

Главный эксперимент~--- честное сравнение \emph{нейросетевых архитектур}
(табличного MLP и последовательностных GRU/LSTM/TCN) с сильными базовыми моделями
(LightGBM, XGBoost, CatBoost) на основной конфигурации, на \emph{одних и тех же}
purged-фолдах walk-forward. Все обучатели реализуют единый интерфейс и оцениваются
через один и тот же модуль метрик; нейросетевая часть исполняется в отдельном
процессе. Принципиально, что \emph{автоматические тесты утечки} применяются и к
нейросетям (вставка метки как фичи~$\to$~метрики близки к идеалу; перемешивание
таргета~$\to$~падение до baseline; сдвиг признаков на горизонт вперёд~$\to$~рост),
а фолды нейросетевого субпроцесса побитово совпадают с фолдами базлайнов. Это
позволяет ответить на центральный вопрос работы: оправдывает ли выразительность
нейронных сетей их склонность к переобучению на высоковолатильном ряде с малым
числом независимых наблюдений.

\section{Рамка постановок и абляции}

Принципиальный эксперимент встроен в рамку, фиксирующую влияние постановки задачи и
определения целевой переменной: небольшой пред-зарегистрированный грид сквозных
экспериментов по оси \emph{типа целевой переменной} (point-to-point, triple-barrier
и регим-нормализованный vol-scaled с несколькими значениями~$k$), внутри каждого из
которых сравниваются три постановки (A,~B,~C, причём C в двух вариантах c1/c2) на
фиксированном обучателе, признаках, фолдах и пороге. Число конфигураций
\emph{сознательно ограничено}: каждый дополнительный эксперимент повышает планку
статистической значимости (через число испытаний в DSR и в оценке вероятности
переобучения). Дополнительно предусмотрены абляции: режим окна (расширяющееся против
скользящего) и взвешивание по давности~--- для оценки чувствительности к
нестационарности; вклад деривативных признаков (с/без funding и базиса); вклад
кросс-активов на коротком окне. Диагностика переобучения включает сопоставление
качества на обучении и вне выборки, важность признаков методом перестановок и
значения Шепли, а также мониторинг распада эджа и дрейфа распределений (PSI/KS) по
ходу walk-forward. Фиксированные пороги, отличные от основного~$2\%$, заведены как
механизм, но их полный экономический свип отнесён к направлениям дальнейшей работы.

\section{Оборудование и воспроизводимость}

Поскольку набор данных невелик, основная вычислительная нагрузка~--- это
\emph{CPU-параллельная} задача пропускной способности (множество мелких независимых
обучений), а не GPU-задача; графический ускоритель необходим лишь нейросетевым
моделям. Для эффективного использования ресурсов вводится единый тред-бюджет
(произведение числа параллельных обучений на число потоков одного обучения не
превышает числа физических ядер), нейросетевая часть выносится в отдельный процесс,
а параллельный подбор гиперпараметров используется для ускорения при сохранении
детерминизма финальной переоценки через зафиксированные наилучшие параметры. Полная
воспроизводимость обеспечивается фиксацией всех источников случайности, фиксированной
датой отсечки данных и сохранением сырого снимка с контрольной суммой.

\section{Результаты}
\label{sec:results}

\noindent\emph{[Результаты будут добавлены после полного вычислительного прогона.]}

\medskip
В данном разделе будут приведены: таблица сравнения \emph{нейросетевых архитектур}
(MLP, GRU, LSTM, TCN) с базовыми моделями (LightGBM, XGBoost, CatBoost) на
идентичных фолдах; сводная таблица сравнения постановок A/B/C по классификационным
(MCC, macro-$F_1$, полнота по классам $\pm 1$), сигнальным (информационный
коэффициент и его стабильность) и экономическим (коэффициент Шарпа, Deflated Sharpe
Ratio) метрикам; распределения метрик по фолдам (диаграммы размаха) вместо одиночных
чисел; наложенные кривые precision-at-coverage; результаты теста Дибольда--Мариано
между моделями и постановками; результаты абляций; а также итоги статистического
контроля переобучения (DSR, PBO)
по всему гриду. Никакие числовые результаты до завершения прогона не приводятся.

% =====================================================================
%  ЗАКЛЮЧЕНИЕ
% =====================================================================
\chapter*{Заключение}
\addcontentsline{toc}{chapter}{Заключение}

\noindent\emph{[Будет написано по результатам исследования.]}

\medskip
В заключении предполагается обобщить итоги сравнения трёх постановок задачи
прогнозирования направления цены биткоина, сформулировать выводы относительно
границ предсказуемости направления на часовом горизонте, соотнести полученные
результаты с поставленными задачами и наметить направления дальнейшей работы
(явная модель режима на основе скрытой марковской модели; полный экономический свип
порога и типа целевой переменной; последовательностные модели на большем объёме
данных; обогащение экономического слоя моделью размера позиции и издержек;
кросс-биржевая и мультиактивная проверка обобщаемости; живая (forward) валидация).

% =====================================================================
%  СПИСОК ЛИТЕРАТУРЫ
% =====================================================================
\begin{thebibliography}{99}
\addcontentsline{toc}{chapter}{Список литературы}

\bibitem{nakamoto2008} Nakamoto~S. Bitcoin: A Peer-to-Peer Electronic Cash System.
2008. URL: \url{https://bitcoin.org/bitcoin.pdf}.

\bibitem{fama1970} Fama~E.~F. Efficient Capital Markets: A Review of Theory and
Empirical Work~// The Journal of Finance. 1970. Vol.~25, no.~2. P.~383--417.

\bibitem{mandelbrot1963} Mandelbrot~B. The Variation of Certain Speculative
Prices~// The Journal of Business. 1963. Vol.~36, no.~4. P.~394--419.

\bibitem{lo1988} Lo~A.~W., MacKinlay~A.~C. Stock Market Prices Do Not Follow Random
Walks: Evidence from a Simple Specification Test~// The Review of Financial Studies.
1988. Vol.~1, no.~1. P.~41--66.

\bibitem{ji2019} Ji~S., Kim~J., Im~H. A Comparative Study of Bitcoin Price
Prediction Using Deep Learning~// Mathematics. 2019. Vol.~7, no.~10. Art.~898.
DOI:~10.3390/math7100898.

\bibitem{livieris2020} Livieris~I.~E., Pintelas~E., Stavroyiannis~S., Pintelas~P.
Ensemble Deep Learning Models for Forecasting Cryptocurrency Time-Series~//
Algorithms. 2020. Vol.~13, no.~5. Art.~121. DOI:~10.3390/a13050121.

\bibitem{ortu2022} Ortu~M., Uras~N., Conversano~C., Bartolucci~S., Destefanis~G.
On Technical Trading and Social Media Indicators for Cryptocurrency Price
Classification through Deep Learning~// Expert Systems with Applications. 2022.
Vol.~198. Art.~116804. DOI:~10.1016/j.eswa.2022.116804.

\bibitem{kim2022} Kim~G., Shin~D.-H., Choi~J.~G., Lim~S. A Deep Learning-Based
Cryptocurrency Price Prediction Model That Uses On-Chain Data~// IEEE Access. 2022.
Vol.~10. P.~56232--56248. DOI:~10.1109/ACCESS.2022.3177888.

\bibitem{dashtaki2025} Mohammadi Dashtaki~S., Hosseini Chagahi~M., Bahadori~A.,
Moshiri~B., Piran~M.~J., Montazeri~A. HSIF: A Transformer-Based Cross-Attention
Framework for Cryptocurrency Trend Forecasting via Multimodal Sentiment--Market
Fusion~// IEEE Access. 2025. Vol.~13. DOI:~10.1109/ACCESS.2025.3605522.

\bibitem{wu2024} Wu~J., Zhang~X., Huang~F., Zhou~H., Chandra~R. Review of Deep
Learning Models for Crypto Price Prediction: Implementation and Evaluation. 2024.
arXiv:2405.11431. URL: \url{https://arxiv.org/abs/2405.11431}.

\bibitem{labbaf2024} Labbaf Khaniki~M.~A., Manthouri~M. Enhancing Price Prediction
in Cryptocurrency Using Transformer Neural Network and Technical Indicators. 2024.
arXiv:2403.03606. URL: \url{https://arxiv.org/abs/2403.03606}.

\bibitem{lopez2018} L\'opez de Prado~M. Advances in Financial Machine Learning.
Hoboken, NJ: John Wiley~\&~Sons, 2018.

\bibitem{bailey2014} Bailey~D.~H., L\'opez de Prado~M. The Deflated Sharpe Ratio:
Correcting for Selection Bias, Backtest Overfitting, and Non-Normality~// The
Journal of Portfolio Management. 2014. Vol.~40, no.~5. P.~94--107.

\bibitem{bailey2017} Bailey~D.~H., Borwein~J., L\'opez de Prado~M., Zhu~Q.~J. The
Probability of Backtest Overfitting~// Journal of Computational Finance. 2017.
Vol.~20, no.~4. P.~39--69.

\bibitem{diebold1995} Diebold~F.~X., Mariano~R.~S. Comparing Predictive Accuracy~//
Journal of Business~\&~Economic Statistics. 1995. Vol.~13, no.~3. P.~253--263.

\bibitem{matthews1975} Matthews~B.~W. Comparison of the Predicted and Observed
Secondary Structure of T4 Phage Lysozyme~// Biochimica et Biophysica Acta (BBA)~---
Protein Structure. 1975. Vol.~405, no.~2. P.~442--451.

\bibitem{ke2017} Ke~G., Meng~Q., Finley~T., Wang~T., Chen~W., Ma~W., Ye~Q.,
Liu~T.-Y. LightGBM: A Highly Efficient Gradient Boosting Decision Tree~//
Advances in Neural Information Processing Systems (NeurIPS). 2017. Vol.~30.
P.~3146--3154.

\bibitem{chen2016} Chen~T., Guestrin~C. XGBoost: A Scalable Tree Boosting System~//
Proceedings of the 22nd ACM SIGKDD International Conference on Knowledge Discovery
and Data Mining. 2016. P.~785--794.

\bibitem{prokhorenkova2018} Prokhorenkova~L., Gusev~G., Vorobev~A., Dorogush~A.~V.,
Gulin~A. CatBoost: Unbiased Boosting with Categorical Features~// Advances in
Neural Information Processing Systems (NeurIPS). 2018. Vol.~31. P.~6638--6648.

\end{thebibliography}

\end{document}
